% !PRINTED 89
\weberpage{89}
%<<CONT

\begin{equation}
 h=\nu+1,\nu+2,\ldots,n,
 \qquad
 k=\nu+1,\nu+2,\ldots,m,
\tag{4}
\end{equation}
and hence
\begin{equation}
        (n-\nu)(m-\nu)
\tag{5}
\end{equation}
in number.

These conditions are, however, actually mutually independent, that is, none follows from the others; for the left-hand sides of (3), by suitable assumptions concerning the coefficients \(a\), can receive an entirely arbitrary value for every combination of indices in the series (4), as is seen by setting all \(a_h^{(i)},a_\mu^{(k)}\), with the exception of \(a_h^{(k)}\), equal to zero.

\begin{center}
{\large \S.\ 26.}

\vspace{0.5em}
{\large\gesperrt{Non-homogeneous Linear Equations.}}
\end{center}

We have simplified the problem of solving linear equations in the preceding considerations, and reduced it to more general laws, to no inconsiderable degree by assuming the equations homogeneous with respect to the unknowns. In applications, however, the unknowns frequently do not occur homogeneously, and although in principle one case is not essentially different from the other, we shall nevertheless consider the case of non-homogeneous equations separately.

We can derive it from the case of homogeneous equations by adding the requirement that one specified unknown shall have the value \(1\). If this unknown occurs among those which the problem leaves arbitrary, no difficulty arises, because we need only put it \(=1\). But if it belongs to those determined by the others, certain further conditions must be satisfied, which say that the value \(1\) is admissible for this unknown.

We shall treat the question here independently and with a somewhat altered notation, but for simplicity restrict ourselves to the most important case, in which the number of unknowns agrees with the number of equations.

Let the following system of equations be solved with respect to the unknowns \(x_1,x_2,\ldots,x_n\):
%CONT>>
